\subsection{Método de Newton}

Consideraremos buscar el cero de nuestra función f(x) a través del uso de la serie de Taylor.

Para esto, tomemos en nuestro algoritmo la función $f(x)$, si desarrollamos en segundo orden la serie de Taylor para el cero, se puede notar que:

\begin{equation}
f(x_0)= f(x) + (x - x_0) \cdot f'(x) + (x - x_0)^2 \cdot \frac{f''(x)}{2} + ...
\end{equation}

Pero por definición, la función evaluada en nuestro punto $x_0$ debe ser cero, así, si aproximamos la serie:

\begin{equation}
0 = f(x) + (x-x_0) \cdot f'(x)
\end{equation}

Podemos despejar así $x_0$, donde:

\begin{equation}
x_0 = x - \frac{f(x)}{f'(x)}
\end{equation}

Esto, nos ofrece un algoritmo que seguir, a modo de poder aproximarse a través de iteraciones hasta superar la tolerancia establecida.

Consideraremos un punto inicial, tal como el valor inicial del intervalo dentro del cual estamos buscando el cero, para este caso, será $x_n$, así, la iteración del algoritmo nos dará un valor cercano al cero, siendo este denotado $x_{n+1}$, así:

\begin{equation}
x_{n+1} = x_n - \frac{f(x_n)}{f'(x_n)}
\end{equation}

Un ejemplo de la implementación de este algoritmo es el siguiente, consideramos una función $f(x) = x^3 -  5x$, y su derivada $f'(x) = 3x^2 - 5$, así:

\begin{lstlisting}
tolerancia = 1e-12

f = lambda x: x**3 -5*x # Se define la función a utilizar
fprima = lambda x: 3*x**2 - 5

def met_newton(f,a,b,tolerancia):
    zeros = []
    int = np.linspace(a,b,10)
    print(int)
    for i in range(int.size-1): 

        if f(int[i])*f(int[i+1]) <= 0:
            a,b = int[i:i+2]
            fa, fprima_a = f(a), fprima(a)
            m = 0
            while (abs(fa/fprima_a )) >= tolerancia:
                b = a - fa/fprima_a
                a = b
                fa, fprima_a = f(a), fprima(a)
                m +=1
            print("Existe un cero",b,"con iteraciones de",m,"veces")


met_newton(f,-3,3,tolerancia)

\end{lstlisting}

Es importante mencionar que el método de Newton es uno de los métodos más rápidos a la hora de encontrar ceros, esto, gracias al uso de la derivada de la función.

A su vez, notar que si la derivada se acerca a cero, la cantidad de iteraciones comenzará a aumentar, puesto que la velocidad con la que avanza el algoritmo comienza a disminuir.

